\chapter*{Avant-propos}
\addcontentsline{toc}{chapter}{Avant-propos}

Ce rapport est rédigé dans le cadre de la validation de la première
année du Master Informatique à l'Université Paris 8 Vincennes -- Saint-Denis,
au sein de l'UFR des Sciences et Technologies du Numérique (STN --
MITSIC).

Il présente un travail de recherche portant sur la détection
d'anomalies dans le trafic des réseaux cellulaires 4G/5G, mené à partir
de données réelles de trafic mobile de la ville de Milan. Ce sujet
s'inscrit dans la continuité des travaux menés sur le framework
DiNATrAX (Maudoux \& Boumerdassi, IEEE ICC 2024), dont l'approche par
encodage de signatures et distance d'édition a servi de point de départ
méthodologique à ce projet.

Au-delà de la mise en \oe uvre technique, ce travail a été l'occasion
d'approfondir des notions de traitement de séries temporelles,
d'algorithmique des chaînes de caractères et d'analyse statistique,
tout en les confrontant à un cas d'usage concret et à des données
réelles de grande échelle.
